\section{The deep-nesting pair}\label{sec:dnp}

The deep-nesting pairs of \cref{sec:prelim-dnp} have two heads of unbounded owner-load as $m$ grows, so they escape
Theorems F and~M, and they remain the sharpest open instance of $\TCP(2)$. This section proves three things about them.
A single turn counter serves every level of both heads at $\log_2q=m+O(\log(1/\theta))$ and cuts the pair at depth
$2^{m+O(1)}/\eps$ (\cref{sec:dnp-counter}); so Conjecture~A$'$ can hold only with $C\ge1$, and the depth question is
exponential against subexponential in~$m$. Fine levels are \emph{blind}: on average over shifts, a class with $g$ levels
below it sees a turn of the other head on at most a fraction $(g+3+\lceil\log_2(g+2)\rceil)/(m+1)$ of its hops, for every
label that changes only at turns (\cref{sec:dnp-blind}). Coarse levels are \emph{visible}: for the block labels that serve one head
cheaply, a fixed fraction of the apex hops of every coarse class of the other head sees a cancellation at full size, an
other-head increment of at most $-Q-1$, on average over shifts (\cref{sec:dnp-visible}). Blindness leaves a weaker form of Conjecture~A$'$, which still gives
exponential depth (\cref{conj:App}).

Throughout, $x_1=U(s_0,m,1)$ and $x_2=x_1(\cdot+\xi)$ for an integer shift~$\xi$, with $s_0>m$, $\Lambda=2^{s_0}$,
$R_i=2^{s_0-i}$, $n_i=m-i+1$ and $P=4(m+1)\Lambda$, as in \cref{sec:prelim-dnp}.

\subsection{Turns and classes}\label{sec:dnp-classes}

Let $B$ be a V of level $k\ge1$ of a head, attached at its base~$b$ to a stroke of direction~$\sigma$, with apex~$\alpha$,
valley~$\nu$ and valley cell $v=b-\sigma R_k$. Its back stroke is a segment from $b$ to~$v$ followed, if $k<m$, by the V
$C_{k+1}$ of level $k+1$ attached at~$v$; for $k+1<l\le m$ let $C_l$ be the V of level~$l$ attached at the end of the segment
of the forward stroke of $C_{l-1}$. By \cite[Lemma~5.4(a)]{FourthPaper2026}, which has a written proof there, $C_l$ is attached at~$v$, has valley cell
$v+\sigma R_l$, stays between $v$ and $v+\sigma R_l$, and ends at time~$\nu$, since $C_{k+1}$ ends the back stroke of~$B$ and
$C_{l+1}$ ends the forward stroke of~$C_l$; its valley time is $\nu-n_lR_l$. We call $C_{k+1},\dots,C_m$ the \emph{chain}
of~$B$, and put $R_{m+1}=0$.

\begin{lemma}[Turns and classes]\label{lem:dnp-classes}
\notformalized
\begin{enumerate}[label=(\alph*)]
\item Each head has $2^{m+3}-6$ turns per period, and the pair has $N_P=2^{m+4}-12$.
\item For $1\le k\le m$, the hops of level~$k$ of a head are exactly the following $2R_k$ hops for each V~$B$ of level~$k$ of
that head, in the notation above: the \emph{apex hops} $(\alpha-d,\alpha+d)$, $1\le d\le R_k$, and the \emph{valley hops} at
the cells $v+\sigma x$, $1\le x\le R_k$, whose later endpoint is $\nu+x$. For $d<R_k$ the only turn of the head in the
half-open interval $(\alpha-d,\alpha+d]$ of an apex hop is~$\alpha$. For $k\le l\le m$ and $R_{l+1}<x<R_l$, the valley hop at
$v+\sigma x$ is
\[
  \bigl(\nu-(m-l)R_l-x,\ \nu+x\bigr),
\]
and the turns of the head in its half-open interval are $\nu$ and the turns of $C_{l+1}$ (only $\nu$ if $l=m$).
\item On the $r$-fold cover each head has exactly $4r\Lambda$ hops of each level $i\in\{0,\dots,m\}$; so every class has
$n_{\mathrm h}/(2(m+1))$ hops, and the hops of~(b) other than those listed with $d<R_k$ or $R_{l+1}<x<R_l$, the
\emph{boundary hops}, are a fraction $(m-k+2)/(2R_k)$ of the class.
\end{enumerate}
\end{lemma}

\begin{informalproof}
(a) The turns are the apex and valley times of the V's, two per V, and two tips per period; a period contains $2^{i+1}$ V's
of level~$i$ (\cref{sec:prelim-dnp}). So each head has $2+\sum_{i=1}^m2^{i+2}=2^{m+3}-6$ turns per period.

(b) The V~$B$ is attached at the end of the segment of a stroke of level $k-1$, which has $R_{k-1}=2R_k$ steps in
direction~$\sigma$ and ends at~$b$ at time~$\alpha$; the back stroke starts with a segment of $R_k$ steps in
direction~$-\sigma$. So for $1\le d\le R_k$ the cell $b-\sigma d$ is visited at $\alpha-d$ and at $\alpha+d$ and not in
between, and for $d<R_k$ the only turn in $(\alpha-d,\alpha+d]$ is~$\alpha$. For the valley hops, the forward stroke of~$B$
starts with a segment from~$v$ to~$b$, so $v+\sigma x$ is visited at $\nu+x$ for $1\le x\le R_k$, and before that time,
after~$\nu$, it is not visited. Let $R_{l+1}<x<R_l$. If $l=k$, the chain stays between $v$ and $v+\sigma R_{k+1}$, so the
last visit before~$\nu$ is by the back segment, at time $\alpha+R_k-x=\nu-(m-k)R_k-x$, using $\nu=\alpha+n_kR_k$. If
$l>k$, the forward segment of~$C_l$ runs from $v+\sigma R_l$ to~$v$ starting at time $\nu-n_lR_l$, and afterwards $C_{l+1}$
stays between $v$ and $v+\sigma R_{l+1}$; so the last visit is at $\nu-n_lR_l+R_l-x=\nu-(m-l)R_l-x$, after the valley
time of~$C_l$. In both cases the turns in the half-open interval are those of $C_{l+1}$ and~$\nu$, because $\nu+x$ lies
before the apex of the V of the forward stroke, at $\nu+R_k$.

The listed hops have level exactly~$k$. The only turn strictly inside an apex hop is~$\alpha$, since the times before~$\alpha$
in the hop belong to the parent's segment and those after it to the back segment. A valley hop at $v+\sigma x$ has its
earlier endpoint at or after~$\alpha$, because the back segment visits its cell at time $\alpha+R_k-x$, and its later
endpoint is $\nu+x\le\nu+R_k$; so its interior lies in the span of~$B$, whose turns have level at least~$k$, and it contains
the turn~$\nu$ of level~$k$. Each listed hop contains exactly one turn of level~$k$, the apex or the valley of~$B$, so it
determines~$B$ and its type; within~$B$, the apex hops are distinguished by~$d$ and the valley hops by their cells. So the
listed hops are distinct, and every head has at least $2^{k+1}\cdot2R_k=4\Lambda$ hops of level~$k$ per period, for
$1\le k\le m$.

(c) A hop has level~$0$ exactly when it contains a tip strictly inside. At a tip~$\tau$ the head is at $\pm\Lambda$, and for
each of the other $2\Lambda$ cells of $[-\Lambda,\Lambda]$ exactly one hop of the head at that cell has $\tau$ in its
half-open interval, hence in its open interval. No hop contains two tips: between consecutive tips the head sweeps all of
$[-\Lambda,\Lambda]$, the second and third level-$0$ strokes of a period sweeping from $\Lambda$ to $-\Lambda$ and the fourth
and first from $-\Lambda$ to $\Lambda$, and a hop at~$\pm\Lambda$ would contain the tip at its own cell. So each head has
exactly $4\Lambda$ hops of level~$0$ per period. With~(b), every level has at least $4\Lambda$ hops per period and head, and
there are exactly $P=4(m+1)\Lambda$ hops per period and head, one ending at every time. Hence every level has exactly
$4\Lambda$, the list of~(b) is complete, and on the $r$-fold cover every class has $4r\Lambda=n_{\mathrm h}/(2(m+1))$ hops.
The boundary hops are $m-k+2$ of the $2R_k$ hops of each V: $d=R_k$ and $x=R_l$ for $k\le l\le m$.
\end{informalproof}

The fourth paper defined service with the normalization $n_{\mathrm h}/(2(m+1))$ per class on the strength of numerical
counts~\cite[Section~5.7]{FourthPaper2026}; \cref{lem:dnp-classes}(c) proves that every class has exactly this size. All
turn times of head~$1$ are multiples of~$R_m$, and those of head~$2$ are congruent to $-\xi$ modulo~$R_m$; so for odd
$\xi$ no time is a turn of both heads.

\subsection{The turn counter}\label{sec:dnp-counter}

\begin{theorem}[Turn counter]\label{thm:counter}
\partformal{TurnCounter.lean}{turnCounter_cost}{part~(1), the hypotheses of the Lean theorem on the cover, and the counts
of \cref{lem:dnp-classes}(a) and \cref{sec:prelim-dnp}, which rest on the fourth paper's Lemma~5.4 and its count of V's.}
Let $\xi$ be any shift, $q\ge3$, and $r$ a multiple of $q/\gcd(q,N_P)$. On the $r$-fold cover let $\mathcal N(t)$, for
$t\in[0,T)$, be the number of turns of either head at times in $[0,t]$, and $G(t)=\mathcal N(t)\bmod q$. Then:
\begin{enumerate}[label=(\arabic*)]
\item every hop $(p,u)$ has error $e=(\kappa-1)\bmod q\le\kappa-1$, where $\kappa\ge1$ is the number of turns of either
head in the lifted interval $(p,u]$;
\item $\desc(G)=rN_P/q$ and $q\cdot\cost(G)\le rN_P(4\Lambda+3)$.
\end{enumerate}
\end{theorem}

\begin{outline}[Outline of the formal proof of the label arithmetic]
The Lean theorem \lean{turnCounter_cost} concerns a finite set of times with a cyclic predecessor bijection and a finite set
of hops arranged in closed chains, as in Lemma~Zq (\lean{MotionDepthCodes.lemma_Zq}). It assumes a label $G$ with values in
$[0,q)$ that advances, modulo~$q$, by $\chi_t\in[0,q)$ at every time~$t$ and by $\kappa_h\ge1$ along every hop~$h$, and it
concludes $q\cdot\cost(G)\le\sum_h\kappa_h+\sum_t\chi_t$ and $q\cdot\desc(G)=\sum_t\chi_t$. Its proof: Lemma~Zq writes
$q\cdot\cost(G)$ as the number of hops plus the sum of the reduced errors plus $\sum_t\chi_t$, and each reduced error
$(G(u)-G(p)-1)\bmod q=(\kappa_h-1)\bmod q$ is at most $\kappa_h-1$.
\end{outline}

\begin{proof}[Proof of~(1), of the hypotheses and of the counts (not formalized)]
The times form the cycle $\Z_T$ under $t\mapsto t-1$, and the hops of each head at each cell form closed chains, each hop
followed by the hop that starts at its later endpoint, since every time is the earlier and the later endpoint of exactly one
hop of each head (\cref{sec:prelim-dnp}). Extend $\mathcal N$ to~$\Z$ by $\mathcal N(t+T)=\mathcal N(t)+rN_P$. Since $r$ is a multiple of $q/\gcd(q,N_P)$, $q$
divides $rN_P$, so $G\equiv\mathcal N\pmod q$ on all of~$\Z$. Hence $G$ advances at each time~$t$ by the number $\chi_t$
of turns at~$t$, at most $2<q$, and along each hop by the number~$\kappa$ of turns in its lifted half-open interval, which
is at least~$1$ because the owner of the hop is a turn strictly inside it (\cref{sec:prelim-walks}). This gives~(1), since
$e\equiv\kappa-1\pmod q$ and $\kappa-1\ge0$. Each turn is counted once in $\sum_t\chi_t$, so $\sum_t\chi_t=rN_P$ by
\cref{lem:dnp-classes}(a). Each turn lies in the half-open intervals of exactly $2(2\Lambda+1)$ hops (\cref{sec:prelim-dnp}),
so $\sum_h\kappa_h=rN_P\cdot2(2\Lambda+1)$. The Lean theorem then gives $\desc(G)=rN_P/q$ and
$q\cdot\cost(G)\le rN_P(4\Lambda+2)+rN_P$.
\end{proof}

\begin{corollary}[The counter serves every level]\label{cor:counter-served}
\notformalized
In the setting of \cref{thm:counter}, let $0<\theta<1$. If $q\ge N_P(1+1/(2\Lambda))/\theta$, for instance if
$q\ge20\cdot2^m/\theta$, then every level is doubly $\theta$-served. If moreover $q\ge32\cdot2^m/\theta$ and
$\theta\le1/3$, then $\cost(G)\le\frac{11}{32}\,\theta\,n_{\mathrm h}/(m+1)\le n_{\mathrm h}/8$.
\end{corollary}

\begin{informalproof}
By \cref{thm:counter}(1), the total error of all hops is at most $\sum_h(\kappa_h-1)<rN_P(4\Lambda+2)$, so every class has
total error below this. It is at most $\theta qn_{\mathrm h}/(2(m+1))=4\theta qr\Lambda$ as soon as
$q\ge N_P(4\Lambda+2)/(4\theta\Lambda)=N_P(1+1/(2\Lambda))/\theta$; and $N_P<2^{m+4}$, $\Lambda\ge2$. By
\cref{thm:counter}(2), $\cost(G)/n_{\mathrm h}\le N_P(4\Lambda+3)/(8(m+1)\Lambda q)$, which for $q\ge32\cdot2^m/\theta$ is
at most $\theta(4\Lambda+3)/(16(m+1)\Lambda)\le\frac{11}{32}\theta/(m+1)$, since $(4\Lambda+3)/\Lambda\le11/2$.
\end{informalproof}

\begin{corollary}[Conjecture A$'$ needs $C\ge1$]\label{cor:C1}
\notformalized
If Conjecture~A$'$ holds with constants $(\theta_0,C,C')$, then $C\ge1$.
\end{corollary}

\begin{informalproof}
Put $\theta=\min(\theta_0,1/3)$ and $q=2^{m+5+\lceil\log_2(1/\theta)\rceil}\ge32\cdot2^m/\theta$. For every~$m$, every
$s_0>m$ and every shift, in particular the one provided by the conjecture, \cref{cor:counter-served} gives a label on a cover of
cost at most $n_{\mathrm h}/8$ that doubly $\theta_0$-serves all $m+1$ levels. Hence
$m+1\le C(m+5+\lceil\log_2(1/\theta)\rceil)+C'$ for all~$m$, and letting $m\to\infty$ gives $C\ge1$.
\end{informalproof}

\begin{corollary}[Cutting the pair with a counter]\label{cor:counter-cut}
\notformalized
For every shift, every $\eps>0$ and every multiple~$T$ of~$P$, the walk DAG of the pair on $[0,T)$ has a cut of at most
$\eps T$ edges and depth at most $\lceil11\cdot2^{m+1}/((m+1)\eps)\rceil-1$. For $m\ge1$ and $\eps\le1$ this is smaller
than the depth $\lceil12/\eps\rceil^{m+1}-1$ of the fourth paper's clock labels.
\end{corollary}

\begin{informalproof}
Let $\mathcal N(t)$ be the number of turns of either head at times in $[0,t]$, $\Psi(t)=\mathcal N(t)\bmod q$, and call $t$
a \emph{wrap time} if $\lfloor\mathcal N(t)/q\rfloor>\lfloor\mathcal N(t-1)/q\rfloor$. A hop $(p,u)$ with no wrap time
in $(p,u]$ is not violated, since then $\Psi(u)-\Psi(p)=\mathcal N(u)-\mathcal N(p)\ge1$ by its owner. Descents occur only
at wrap times. There are at most $\mathcal N(T-1)/q\le(T/P)N_P/q$ wrap times, and each lies in the half-open intervals of
at most $2(2\Lambda+1)$ hops. So $\cost(\Psi)\le(T/P)N_P(4\Lambda+3)/q$, which is at most $\eps T$ for
$q\ge N_P(4\Lambda+3)/(4(m+1)\Lambda\eps)$; this holds for $q=\lceil11\cdot2^{m+1}/((m+1)\eps)\rceil$, because
$N_P<2^{m+4}$ and $(4\Lambda+3)/(4\Lambda)\le11/8$. \Cref{lem:labels} gives depth at most $q-1$. For the comparison,
$11\cdot2^{m+1}/(m+1)\le12^{m+1}/6$ for $m\ge1$ (at $m=1$ it reads $22\le24$, and from $m$ to $m+1$ the left side less
than doubles while the right side grows twelvefold), so with $x=12^{m+1}/\eps\ge144$ we get
$\lceil11\cdot2^{m+1}/((m+1)\eps)\rceil\le\lceil x/6\rceil<x\le(12/\eps)^{m+1}\le\lceil12/\eps\rceil^{m+1}$ for
$\eps\le1$.
\end{informalproof}

\paragraph{Why they matter.} The counter shows what Conjecture~A$'$ must look like: all $m+1$ levels are doubly served
at $\log_2q=m+O(\log(1/\theta))$, within the cost budget, at every shift. So the conjecture, if true, is sharp up to the
constant~$C$. The least depth $D(\eps)$ of a cut of $\eps T$ edges of the walk DAG satisfies
$\log_2D(\eps)\le m+\log_2(1/\eps)+O(1)$ unconditionally, by \cref{cor:counter-cut}, while Conjecture~A$'$ would give $\log_2D(\eps)\ge c(m+1)-c'$ for small~$\eps$ (\cref{cor:wrap}), with $c\le1$ by the counter. So the open question is
exponential against subexponential depth in the nesting depth~$m$, and the base of an exponential is at most~$2$, not the
$c/\eps$ of the clock labels.

\begin{remark}[Numerical: the evidence for Conjecture A$'$]\label{rem:counter-evidence}
The fourth paper quotes, as evidence for Conjecture~A$'$, the numbers $1$, $1$, $2$ and $2$ or~$3$ of doubly
$0.1$-served levels of optimized labels at $m=3$, $s_0=6$ and $q=8,16,32,64$~\cite[Remark~5.26]{FourthPaper2026}. On the
walk DAG of eight periods at the golden shift, the turn counter doubly serves $1$, $1$, $2$ and $3$ levels at the same~$q$,
and at $m=6$, $s_0=10$ it serves $\log_2q-3$ levels for $\log_2q=7,8,9$. The quoted evidence is therefore essentially
what the counter does; it is consistent with Conjecture~A$'$ with $C\approx1$, but it cannot discriminate. (Not all of
these labels meet the cost budget; the comparison is of served counts against~$q$.)
\end{remark}

The counter is a special case of a general criterion.

\begin{proposition}[The counter criterion]\label{prop:counter-criterion}
\notformalized
Let $x_1,\dots,x_k$ be walks of length~$T$ without pauses, with $\mathcal N$ turns in all and at most $H$ hops spanning any
time. Then for every $\eps>0$ the walk DAG has a cut of at most $\eps T$ edges and depth at most
$\max(0,\lceil\mathcal N(H+1)/(\eps T)\rceil-1)$. Hence TCP holds for every family with $(\mathcal N/T)\cdot H$ bounded.
\end{proposition}

\begin{informalproof}
Let $\mathcal N(t)$ count the turns of all walks at times in $[0,t]$, each walk separately, and $\Psi(t)=\mathcal N(t)\bmod q$
with $q=\max(1,\lceil\mathcal N(H+1)/(\eps T)\rceil)$. As in \cref{cor:counter-cut}, a hop with no wrap time in its
half-open interval contains its owner, a turn, and is not violated; there are at most $\lfloor\mathcal N/q\rfloor$ wrap
times, each a descent at most once and inside at most $H$ hops. So the cost is at most
$\lfloor\mathcal N/q\rfloor(H+1)\le\eps T$, and the depth is at most $q-1$.
\end{informalproof}

The criterion is incomparable with the fourth paper's Theorem~B: two zigzags of fixed width whose centres drift
monotonically satisfy it but are not home-based. It fails on the deep-nesting pairs as $m\to\infty$, where
$(\mathcal N/T)\cdot H\approx2^{m+4}(4\Lambda+2)/P\approx2^{m+4}/(m+1)$.

The counter also pins down the capacity trade-off of \cref{cor:capacity} on the deep-nesting family. For two walks write
$\mathrm{Cap}(D,\delta,\eps)$ for the statement: for all pairs of walks of length~$T$ there are a set~$R$ of pairs with
$\card R\le\delta T$ and a cut of $G-R$ with at most $\eps T$ edges and depth at most~$D$.

\begin{proposition}[Capacity on the deep-nesting family]\label{prop:dnp-capacity}
\notformalized
\begin{enumerate}[label=(\alph*)]
\item Let $0<\eps<1/2$, $\delta(\eps)=2/\log_2(1/\eps)$ and $D(\eps)=\lceil11/\eps^2\rceil-1$. For every~$m$, every
$s_0>m$, every shift and every multiple $T$ of~$P$, the pair on $[0,T)$ has a set~$R$ of pairs with
$\card R\le\delta(\eps)T$ and a cut of $G-R$ with at most $\eps T$ edges and depth at most $D(\eps)$.
\item Suppose there are $\eps_0>0$, an integer $N_0\ge1$ and a function~$D_*$ such that for every $m\ge0$ there are $s_0>m$
and an odd shift for which the walk DAG of the pair on $[0,N_0P)$ has no cut of at most $\eps_0N_0P$ edges and depth less
than $D_*(m)$ (a walk-DAG form of Conjecture~A$'$, which follows from Conjecture~A$'$ by \cref{cor:wrap}). If $\mathrm{Cap}(D(\eps),\delta(\eps),\eps)$ holds for all~$\eps$,
then $D(\eps)\ge D_*(m)$ whenever $\eps\le\eps_0/2$ and $4(m+1)N_0\delta(\eps)\le\eps_0$. In particular, if
$D_*(m)\ge2^{c(m+1)-c'}$, then $\delta(\eps)\ge c\eps_0/(4N_0(\log_2D(\eps)+c+c'))$ whenever $\eps\le\eps_0/2$ and
$\delta(\eps)\le\eps_0/(4N_0)$.
\end{enumerate}
\end{proposition}

\begin{informalproof}
(a) If $m+1\ge2/\delta(\eps)$, let $R$ be the set of all pairs visited; then $\card R=2(2\Lambda+1)$ and
$\card R/T\le(4\Lambda+2)/(4(m+1)\Lambda)\le(5/4)/(m+1)\le\frac58\delta(\eps)$, every hop is at a pair of~$R$, and the empty
cut of $G-R$ has depth~$0$. If $m+1<2/\delta(\eps)=\log_2(1/\eps)$, let $R=\emptyset$; \cref{cor:counter-cut} gives depth at
most $\lceil11\cdot2^{m+1}/((m+1)\eps)\rceil-1\le\lceil11/\eps^2\rceil-1$, since $2^{m+1}<1/\eps$.

(b) Take the instance for~$m$ and $T=N_0P$, and let $R$ and a cut~$F$ of $G-R$ be given by the capacity statement. Each pair
is visited at most $2(m+1)N_0$ times in $[0,T)$, so it carries fewer than $2(m+1)N_0$ hops. Let $F_0$ be the set of hop edges
of~$G$ all of whose hops lie at pairs of~$R$; then $\card{F_0}\le2(m+1)N_0\card R\le\eps_0T/2$, and every edge of~$G$
outside~$F_0$ is an edge of $G-R$. So $F\cup F_0$ is a cut of~$G$ with at most $\eps T+\eps_0T/2\le\eps_0T$ edges and depth
at most $D(\eps)$, and the hypothesis gives $D(\eps)\ge D_*(m)$. For the second statement put $y=\eps_0/(4N_0\delta(\eps))\ge1$
and $m+1=\lfloor y\rfloor$; then $\log_2D(\eps)\ge c\lfloor y\rfloor-c'\ge c(y-1)-c'$, that is,
$y\le(\log_2D(\eps)+c+c')/c$.
\end{informalproof}

So on the deep-nesting family, under the walk-DAG form of Conjecture~A$'$ assumed in~(b), and so under Conjecture~A$'$
itself (\cref{cor:wrap}), the scratch must have size about
$T/\log D$ when the rest is cut at depth~$D$, and (a) achieves this: there $\log_2D(\eps)\approx2\log_2(1/\eps)+3.5$ and $\delta(\eps)\approx4/\log_2D(\eps)$.

\subsection{Labels that change only at turns}\label{sec:dnp-turnlabels}

Let $\xi$ be odd. A label~$G$ on the cover, or on the walk DAG of a window, is \emph{turn-supported} if there are integers
$J_h(\tau)$ at the turn times~$\tau$ of each head~$h$ such that $G(t)-G(t-1)\equiv J_h(t)\pmod q$ when $t$ is a turn of
head~$h$, and $G(t)\equiv G(t-1)$ when $t$ is not a turn. This class contains the counter ($J\equiv1$), the fourth paper's
clock labels, and every function of the vector of turn counts. The \emph{potential} of head~$h$ is
$\Phi_h(t)=\sum J_h(\tau)$ over the turns $\tau\le t$ of head~$h$ (on a window, or on~$\Z$ for periodic jumps, normalized
by $\Phi_h(0)=0$), so $G\equiv\Phi_1+\Phi_2+\mathrm{const}\pmod q$. For a hop of head~$h'$, the \emph{other-head
increment} is $F=\Phi_h(u)-\Phi_h(p)$, where $h\ne h'$.

For a V of head~$h$ let $M(V)$ be the sum of the jumps of head~$h$ at the turns in the span of~$V$, with $M(\emptyset)=0$;
let $V^\sigb$ and $V^\sigf$ be the V's attached at the ends of the segments of its back and forward strokes (none if $V$ has
level~$m$), and put
\[
  A_V=J_h(\alpha_V)-1,\qquad E_V=J_h(\nu_V)+M(V^\sigb)-1,\qquad c_\sigb(V)=1+A_V,\qquad c_\sigf(V)=2+A_V+E_V .
\]

\begin{lemma}[Path sums]\label{lem:pathsums}
\notformalized
Let $t$ be a time in a level-$0$ stroke~$S_0$ of head~$h$, and let $V_1\supset V_2\supset\dots\supset V_j$ be the V's whose
spans contain~$t$, with $V_i$ of level~$i$ and $V_1$ the V of~$S_0$. Then
\[
  \Phi_h(t)-\Phi_h(\alpha_{V_1}-1)=\sum_{i=1}^jc_{\delta_i}(V_i),
\]
where $\delta_i=\sigb$ if $t$ lies in the back stroke of~$V_i$ and $\delta_i=\sigf$ otherwise; for $t$ in the segment of~$S_0$,
$j=0$. Moreover $M(V)=c_\sigf(V)+M(V^\sigf)$.
\end{lemma}

\begin{informalproof}
By induction on the depth inside one V. The apex $\alpha_V$ is the first time of~$V$. For~$t$ in the segment of the back stroke,
the only turn in $[\alpha_V,t]$ is $\alpha_V$, so the increment over $(\alpha_V-1,t]$ is $J_h(\alpha_V)=c_\sigb(V)$; for $t$
in $V^\sigb$ it is $c_\sigb(V)$ plus the increment from $\alpha_{V^\sigb}-1$ to~$t$. For $t$ in the segment of the forward
stroke the increment is $J_h(\alpha_V)+M(V^\sigb)+J_h(\nu_V)=c_\sigf(V)$, by the definition of $E_V$, and for $t$ in
$V^\sigf$ it is $c_\sigf(V)$ plus the increment from $\alpha_{V^\sigf}-1$ to~$t$. At level~$m$ there are no inner V's and
the same computation applies with $M(\emptyset)=0$. The mass identity is the case where $t$ is the last time of~$V$. The
segment of~$S_0$ contains no turn after its first time.
\end{informalproof}

\begin{lemma}[Band decomposition]\label{lem:band}
\notformalized
Let $G$ be turn-supported, and let $B$ be a V of head~$h'$ of level $k\ge1$, with chain $C_{k+1},\dots,C_m$.
\begin{enumerate}[label=(\alph*)]
\item For $1\le d<R_k$ the apex hop $(\alpha-d,\alpha+d)$ has error $(A_B+F)\bmod q$.
\item For $k\le l\le m$ and $R_{l+1}<x<R_l$, the valley hop at $v+\sigma x$ has error
\[
  \Bigl(E_B-\sum_{i=k+1}^{l}\bigl(A_{C_i}+E_{C_i}+2\bigr)+F\Bigr)\bmod q .
\]
\end{enumerate}
Here $F$ is the other-head increment of the hop.
\end{lemma}

\begin{informalproof}
The error is $(G(u)-G(p)-1)\bmod q$, and $G(u)-G(p)\equiv\Phi_{h'}(u)-\Phi_{h'}(p)+F$. (a) By \cref{lem:dnp-classes}(b) the
only turn of head~$h'$ in $(\alpha-d,\alpha+d]$ is~$\alpha$, with jump $1+A_B$. (b) By \cref{lem:dnp-classes}(b) the
increment of $\Phi_{h'}$ is $M(C_{l+1})+J_{h'}(\nu)$, and $J_{h'}(\nu)=E_B+1-M(C_{k+1})$. Since $C_{i+1}=C_i^\sigf$, the mass
identity of \cref{lem:pathsums} gives $M(C_{k+1})-M(C_{l+1})=\sum_{i=k+1}^l(A_{C_i}+E_{C_i}+2)$.
\end{informalproof}

When $F=0$ on the hops concerned (the \emph{single-head regime}), $A_B$ is the apex error of~$B$ and $c_\sigb-1$,
$c_\sigf-c_\sigb-1$ are errors of $B$'s own hops. So a V whose own errors lie in $[0,E]$ contributes at least $1$ and~$2$ to
the path sums of \cref{lem:pathsums}: the potential of a head that serves many of its levels rises along every path, and it
can come back down only through V's with an apex or band error near~$q$. Serving many levels cheaply needs cancellation,
and the next labels cancel in blocks.

\begin{definition}[Block labels]\label{def:block}
Let $a\ge1$ and $K\ge0$ be integers. A level $l\in[1,m]$ is a \emph{canceller} if $m-l\equiv a\pmod{a+1}$ and
\emph{served} otherwise. Define integers $A_l,E_l$ from level~$m$ upwards: for a served level~$l$, $A_l=K$ and
$E_l=K+\sum_{i=l+1}^{c-1}(A_i+E_i+2)$, where $c$ is the first canceller below~$l$, or $m+1$ if there is none; for a
canceller~$c$, $A_c=0$ and $E_c=-Q_c-2$ with $Q_c=\sum_{i=c+1}^{c'-1}(A_i+E_i+2)$, where $c'$ is the next canceller
below~$c$, or $m+1$. The \emph{block label} $B(a,K)$ on head~$h$ has the jumps $J_h(\alpha_V)=1+A_{\lev V}$,
$J_h(\nu_V)=E_{\lev V}+1-M(V^\sigb)$, computed from level~$m$ upwards, and $J_h=1$ at the tips; so $A_V=A_{\lev V}$ and
$E_V=E_{\lev V}$.
\end{definition}

\begin{lemma}[Block masses]\label{lem:block}
\notformalized
In $B(a,K)$, a served level with $j$ served levels below it in its block has $A+E+2=2^j(2K+2)$. Every canceller has $a$
served levels directly below it, and $Q_c=Q:=(2^a-1)(2K+2)$. Hence $c_\sigb(V)=1+K$ and $c_\sigf(V)=2^j(2K+2)$ at a served
V, and $c_\sigb(V)=1$ and $c_\sigf(V)=-Q$ at a canceller V: the path sums depend only on the level. In the single-head
regime, and if $K+Q<q$, every apex hop with $d<R_k$ of a served V has error~$K$ and every valley hop of
\cref{lem:band}(b) of a served V has error in $[K,K+Q]$.
\end{lemma}

\begin{informalproof}
Write $T(j)$ for $A+E+2$ at a served level with $j$ served levels below it in its block. Then $T(0)=2K+2$ and
$T(j)=2K+2+\sum_{j'<j}T(j')$, so $T(j)=2^j(2K+2)$. A canceller~$c$ satisfies $m-c\equiv a\pmod{a+1}$ and $m-c\ge0$, so
$m-c\ge a$ and the levels $c+1,\dots,c+a$ are served and below~$c$ in one block; hence
$Q_c=\sum_{j<a}T(j)=Q$. The values of $c_\sigb$ and $c_\sigf$ follow. For the errors, write
$S(i_1,i_2)=\sum_{i=i_1}^{i_2}(A_i+E_i+2)$. A served level~$k$ has $E_k=K+S(k+1,c-1)$ with $c$ the first canceller below~$k$.
For $k\le l<c$ the band error of \cref{lem:band}(b) is $K+S(l+1,c-1)\in[K,K+Q)$. For $c\le l<c'$, with $c'$ the next
canceller, it is $K-S(c,l)=K+Q-S(c+1,l)=K+S(l+1,c'-1)\in[K,K+Q]$, because $A_c+E_c+2=-Q$ and $S(c+1,c'-1)=Q$; and the
pattern repeats below~$c'$. All values lie in $[0,q)$, so no reduction occurs.
\end{informalproof}

So, in the single-head regime, one head carrying $B(a,K)$ serves the $a/(a+1)$ fraction of its levels that are served,
with errors in $[K,K+Q]$, and its violated own hops are essentially the valley hops of the canceller classes. Heuristically,
its cost per own hop is then about $1/(2(a+1))$, which for $a=4$ is $1/10$, inside the budget $n_{\mathrm h}/8$ for
every~$m$ at fixed~$q$ (see also \cref{rem:single-head}). So no argument that uses only one head's own structure can be
expected to prove Conjecture~A$'$. At the aligned shift the second head inherits the first head's errors.

\begin{remark}[Numerical: the single-head cost of block labels]\label{rem:single-head}
Exact evaluations on covers gave costs per own hop of $0.05$ to $0.09$ for block labels with $a=4$, $q=256$ and
$4\le m\le10$.
\end{remark}

\begin{proposition}[The aligned shift]\label{prop:aligned}
\notformalized
Let $\xi=1$, put the same jumps on both heads, $J_2(t)=J_1(t+1)$, and for a hop $(a',b')$ of head~$1$ write
$I(a',b')=\Phi_1(b')-\Phi_1(a')$. The hops of head~$2$ are the pairs $(a'-1,b'-1)$ for the hops $(a',b')$ of head~$1$, and
\begin{align*}
  G(b'-1)-G(a'-1)&\equiv2I(a',b')-J_1(b')[b'\text{ a turn}]+J_1(a')[a'\text{ a turn}],\\
  G(b')-G(a')&\equiv2I(a',b')+J_1(b'+1)[b'+1\text{ a turn}]-J_1(a'+1)[a'+1\text{ a turn}].
\end{align*}
So every hop whose shifted endpoints are not turns has error $(1+2e)\bmod q$, where $e=(I-1)\bmod q$ is the single-head
error of the corresponding hop of head~$1$; each turn affects at most four hops.
\end{proposition}

\begin{informalproof}
Head~$2$ is at a cell at time~$t$ exactly when head~$1$ is there at time $t+1$, so its hops are the shifted hops of head~$1$
and its turns are the turns of head~$1$ minus one. Hence $\Phi_2(t)=\Phi_1(t+1)$ and $G(t)\equiv\Phi_1(t)+\Phi_1(t+1)$ up
to a constant, and both identities follow by expanding $\Phi_1(t\pm1)-\Phi_1(t)$. Each time is the earlier and the later
endpoint of exactly one hop of each head.
\end{informalproof}

\begin{remark}[Numerical: Conjecture A$'$ needs a misaligned shift]\label{rem:aligned}
With $B(4,0)$ on both heads at the aligned shift and $q=256$, exact evaluations on covers gave costs $0.080$, $0.0996$ and
$0.0935$ times~$n_{\mathrm h}$ at $m=6$, $9$ and~$10$, inside the budget $n_{\mathrm h}/8$, with $6$ of $7$, $8$ of $10$ and $9$
of $11$ levels doubly $0.1$-served; \cref{prop:aligned} explains this, since at the aligned shift both heads see their
single-head errors, doubled, on all but a vanishing fraction of hops. This indicates that Conjecture~A$'$ cannot hold with the
aligned shift, so that its choice of a misaligned shift is essential. At misaligned shifts
(the golden shift, and a shift whose dyadic phases alternate near $1/3$ and $2/3$ at every level) the same labels
collapse to the finest few levels at costs $0.22$ to $0.24$ times~$n_{\mathrm h}$.
\end{remark}

\subsection{Blindness of fine classes}\label{sec:dnp-blind}

\begin{lemma}[Turn-free runs]\label{lem:runs}
\notformalized
Let $m\ge1$. Over one period, the gaps between consecutive turns of~$U$ have the following lengths: for each level $1\le i\le m$,
$2^{i+2}$ gaps of length~$R_i$, except that two of the gaps of level~$m$ have length $R_m+R_0$ instead; and two gaps of
length~$R_0$.
\end{lemma}

\begin{informalproof}
We find the gap that follows each turn. After an apex of level $i<m$ comes the back segment, $R_i$ steps, and then the apex
of the V attached at its end; after an apex of level~$m$, a plain back stroke of $R_m$ steps and then the valley. After a
valley of level $i<m$ comes the forward segment, $R_i$ steps, and then an apex. After a tip comes the segment of a level-$0$
stroke, $R_0$ steps, and then an apex. After a valley~$\nu$ of level~$m$, the forward stroke takes $R_m$ steps and the V
ends at $\nu+R_m$, together with every stroke of which it is the last part. If some enclosing stroke is a back stroke, the
deepest such ends at $\nu+R_m$ with the valley of its V, a turn. If not, the level-$0$ stroke ends at $\nu+R_m$; this happens
for one level-$m$ V per level-$0$ stroke. If the next level-$0$ stroke is the second or fourth of its period, $\nu+R_m$ is a
tip; otherwise the directions continue (from the second to the third stroke, and from the fourth to the first), the next
segment adds $R_0$ steps, and the next turn is at $\nu+R_m+R_0$; this happens twice per period. Each level~$i$ has
$2^{i+2}$ turns per period, and the lengths add up to $4m\Lambda+4\Lambda=P$.
\end{informalproof}

\begin{proposition}[Blindness]\label{prop:blind}
\notformalized
Let $m\ge1$, let $G$ be turn-supported, let $(p,u)$ be a hop of head~$h'$ of length $\ell=u-p$, and let the shift~$\xi$ be uniformly
random in~$\Z_P$. If $(p,u]$ contains no turn of the other head, then $F=0$. The probability of this event is
\[
  \frac1P\sum_{\text{gaps }L}(L-\ell)^+\;\ge\;\frac1{m+1}\sum_{i=0}^m\Bigl(1-\frac\ell{R_i}\Bigr)^+,
\]
the sum over the gaps of \cref{lem:runs}. For $\xi$ uniform over the odd residues the bound decreases by at most
$2/((m+1)R_m)$.
\end{proposition}

\begin{informalproof}
$F$ is the sum of the other head's jumps over the window, an empty sum. In the other head's own time, where it is~$U$, the
window $(p,u]$ is a window $(a',a'+\ell]$ with $a'$ uniform in~$\Z_P$: a hop of head~$2$, given as a hop $(a,b)$ of~$U$ in the
own time of head~$2$, occupies the window $(a-\xi,b-\xi]$ of head~$1$; and for a hop $(a,b)$ of head~$1$, the turns of head~$2$
are those of~$U$ shifted by~$-\xi$, so the window is $(a+\xi,b+\xi]$ in the time of~$U$. This window contains no turn exactly when
$\tau_j\le a'$ and $a'+\ell<\tau_{j+1}$ for consecutive turns $\tau_j<\tau_{j+1}$, which gives $(L_j-\ell)^+$ values
of~$a'$ for the gap of length $L_j=\tau_{j+1}-\tau_j$. Insert \cref{lem:runs}, using
$(x+y-\ell)^+\ge(x-\ell)^++(y-\ell)^+$ for $x,y,\ell\ge0$ for the two long gaps, and $2^{i+2}/P=1/((m+1)R_i)$,
$4/P=1/((m+1)R_0)$. For odd shifts, each interval of admissible values of~$a'$ contains at least $((L_j-\ell)^+-1)/2$
values of each parity, there are fewer than $2^{m+3}$ gaps, and $2^{m+3}/P=2/((m+1)R_m)$.
\end{informalproof}

\begin{corollary}[Fine classes are blind]\label{cor:blind}
\notformalized
Let $1\le k\le m$, $g=m-k$, and let $\xi$ be uniform in~$\Z_P$. On average over~$\xi$, at least a fraction
\[
  \frac{k-2-\lceil\log_2(g+2)\rceil}{m+1}
\]
of the hops of each class $(h',k)$ contain no turn of the other head. Consequently, for every turn-supported label, the
expected fraction of the hops of the class with $F\ne0$ is at most
\[
  \frac{g+3+\lceil\log_2(g+2)\rceil}{m+1},
\]
which is $O(\log m/m)$ for the classes with $g=O(\log m)$.
\end{corollary}

\begin{informalproof}
By \cref{lem:dnp-classes}(b), the hops of the class are, for each V of level~$k$, its apex hops $(\alpha-d,\alpha+d)$,
$1\le d\le R_k$, and its valley hops at the cells $v+\sigma x$, $1\le x\le R_k$; the boundary hops among them obey the
same length bounds as the others. An apex hop has length $2d\le2R_k=R_{k-1}$, so by \cref{prop:blind} it is blind with
probability at least $\frac1{m+1}\sum_{i=0}^{k-1}(1-2^{i-k+1})\ge\frac{k-2}{m+1}$. For $k\le l\le m$ and
$R_{l+1}<x\le R_l$, the valley hop at $v+\sigma x$ is $(\nu-(m-l)R_l-x,\nu+x)$: for $x<R_l$ by \cref{lem:dnp-classes}(b),
and for the boundary hop $x=R_l$ because the cell $v+\sigma R_l$ is visited at the valley time $\nu-n_lR_l$ of~$C_l$ (at the
apex $\alpha=\nu-n_kR_k$ if $l=k$), left at the next step, and not visited again before $\nu+R_l$. So its length is at most
$(m-l+2)R_l\le R_{l-a}$, where $a=\lceil\log_2(m-l+2)\rceil$; if $l\ge a$ it is blind with probability at least
$\frac1{m+1}\sum_{i=0}^{l-a}(1-2^{i-l+a})\ge\frac{l-a-1}{m+1}\ge\frac{k-1-\lceil\log_2(g+2)\rceil}{m+1}$, using $l\ge k$
and $m-l\le g$; if $l<a$ the claimed bound is negative. The second statement follows, because $F\ne0$ needs a turn of the
other head in the window.
\end{informalproof}

\paragraph{Why it matters.} By Markov's inequality, for each class with $g=O(\log m)$ the set of shifts at which a fixed
fraction $c>0$ of the class has $F\ne0$ has density $O(\log m/(cm))$: for all large~$m$, at all but a vanishing density of
shifts, these classes see the other head on less than a fraction~$c$ of their hops, for every label that changes only at
turns. On their blind hops the error is the single-head error, so a head can serve them as if alone, for instance by a block
label. Heuristically, then, at typical shifts no argument that charges errors to the other head's increments can bound the
number of doubly served fine levels by $C\log_2q+C'+O(\log m)$, and if Conjecture~A$'$ holds at a typical shift, it holds
because serving the coarse levels is expensive. This is not a proof: the corollary bounds an average over shifts, and the
conjecture may choose an atypical shift. The measurements of \cref{rem:blind-numerics} agree with this picture.

\begin{remark}[Numerical: block pairs at a misaligned shift]\label{rem:blind-numerics}
With $B(4,0)$ on both heads at a misaligned shift and $q=256$, the number of doubly $0.1$-served levels was
$4,4,4,5,5,8,8$ at $m=10,15,20,25,30,40,50$, all in the blind fine region, at costs of $0.23$ to $0.29$ times
$n_{\mathrm h}$, above the budget; the value at $m=10$ is exact on a cover, the others are sampled.
\end{remark}

\subsection{Visibility of coarse classes}\label{sec:dnp-visible}

The coarse classes see the other head's cancellations. In the other head's own time an apex hop of half-width~$d$ is a
window $(c_0-d,c_0+d]$ around a centre~$c_0$, and its other-head increment is $F(d)=\Phi_h(c_0+d)-\Phi_h(c_0-d)$.

\begin{lemma}[Straddle involution]\label{lem:straddle}
\notformalized
Let head~$h$ carry $B(a,K)$, let $c_0$ be an integer, $k\ge1$, and let $V$ be a V of head~$h$ of a canceller level~$\kappa$,
with apex~$\alpha_V$, valley~$\nu_V$ and stroke length $\Pi_\kappa=\nu_V-\alpha_V=n_\kappa R_\kappa$. Let $\mathcal D$ be the
set of $d\in[1,R_k-1]$ with $c_0-d\in[\alpha_V,\nu_V)$, $c_0+d\in[\nu_V,\nu_V+\Pi_\kappa)$ and $\Pi_\kappa-d\in[1,R_k-1]$.
Then $d\mapsto\Pi_\kappa-d$ is an involution of~$\mathcal D$, and $F(d)+F(\Pi_\kappa-d)=-2Q-2$ for every $d\in\mathcal D$.
Hence at least half of the elements of~$\mathcal D$ have $F(d)\le-Q-1$.
\end{lemma}

\begin{informalproof}
For $d\in\mathcal D$ put $o_p=c_0-d-\alpha_V$ and $o_u=c_0+d-\nu_V$, both in $[0,\Pi_\kappa)$. For $d'=\Pi_\kappa-d$ we get
$c_0-d'-\alpha_V=o_u$ and $c_0+d'-\nu_V=o_p$, so $d'\in\mathcal D$ and the map swaps $o_p$ and~$o_u$. By \cref{lem:pathsums},
for a time~$x$ in the back stroke of~$V$, $\Phi_h(x)=C_V+c_\sigb(V)+Y(x-\alpha_V)$, and for $x$ in its forward stroke,
$\Phi_h(x)=C_V+c_\sigf(V)+Y(x-\nu_V)$. Here $C_V$ collects $\Phi_h(\alpha_{V_1}-1)$ and the contributions of the V's
containing~$V$, the same for all times in~$V$, and $Y(o)$ is the sum of $c_\delta(W)$ over the V's~$W$ inside a stroke of
level~$\kappa$ that contain its time at offset~$o$; by \cref{lem:block} these contributions depend only on levels, so $Y$ is
the same function for both strokes of~$V$. Hence $F(d)=c_\sigf(V)-c_\sigb(V)+Y(o_u)-Y(o_p)$ and
$F(\Pi_\kappa-d)=c_\sigf(V)-c_\sigb(V)+Y(o_p)-Y(o_u)$, and their sum is $2(c_\sigf(V)-c_\sigb(V))=2(-Q-1)$ by
\cref{lem:block}.
\end{informalproof}

\begin{theorem}[Theorem V: visibility of coarse classes on average over shifts]\label{thm:V}
\notformalized
Let head~$h$ carry $B(a,K)$ and let $\xi$ be uniform in~$\Z_P$. For every class $(h',k)$ of the other head with
$1\le k\le m$ and $g=m-k\ge\lceil\log_2(g+1)\rceil+a$, the expected fraction of its apex hops with $d<R_k$ that have
$F\le-Q-1$ is at least
\[
  \frac{g-\lceil\log_2(g+1)\rceil-2a}{2^{a+4}(m+1)}.
\]
\end{theorem}

\begin{informalproof}
Let $j_1=\min\{j\ge0:g+1-j\le2^j\}$; then $1\le j_1\le\lceil\log_2(g+1)\rceil$, since $g\ge1$. Let $\kappa=k+j$ be the first
canceller level with $j\ge j_1$; cancellers recur every $a+1$ levels, so $j\le j_1+a$, and $\kappa\le m$ by the hypothesis
on~$g$. Since $j\ge j_1$, $\Pi_\kappa=n_\kappa R_\kappa=(g+1-j)R_k2^{-j}\le R_k$, and $\Pi_\kappa$ is even.

Fix a V of level~$k$ of head~$h'$ and its apex hops. As $\xi$ is uniform, the centre~$c_0$ of their windows in the own time of
head~$h$ is uniform modulo~$P$. A window of half-width~$d$ \emph{straddles} a V of level~$\kappa$ of head~$h$ if its left end
lies in the back stroke and its right end in the forward stroke; then $d<\Pi_\kappa$, and $c_0$ lies in the span of that V,
so for fixed~$c_0$ at most one V of level~$\kappa$ is straddled, and for $d\le\Pi_\kappa-1$ the partner
$\Pi_\kappa-d$ lies in $[1,R_k-1]$. So, for fixed~$c_0$, the straddling half-widths form the set~$\mathcal D$ of
\cref{lem:straddle} for the straddled V. For a fixed~$d$, the centres that straddle a given V of level~$\kappa$ form an interval
of length $2\min(d,\Pi_\kappa-d)$, and there are $2^{\kappa+1}$ such V's per period; so
\[
  \mathbb E\card{\mathcal D}=\sum_{d=1}^{\Pi_\kappa-1}\frac{2^{\kappa+2}\min(d,\Pi_\kappa-d)}P
  =\frac{2^{\kappa+2}}P\cdot\frac{\Pi_\kappa^2}4=\frac{2^\kappa\Pi_\kappa^2}P .
\]
By \cref{lem:straddle}, at least $\card{\mathcal D}/2$ of the $R_k-1$ apex hops with $d<R_k$ have $F\le-Q-1$, so the expected
fraction is at least
\[
  \frac{2^{\kappa-1}\Pi_\kappa^2}{PR_k}=\frac{n_\kappa^2R_\kappa}{8(m+1)R_k}=\frac{(g+1-j)^22^{-j}}{8(m+1)},
\]
using $2^\kappa=\Lambda/R_\kappa$ and $P=4(m+1)\Lambda$. By the minimality of~$j_1$, $g+2-j_1>2^{j_1-1}$, so with
$G_1=g+2-j_1$ we have $2^j\le2^{j_1+a}<2^{a+1}G_1$ and $(g+1-j)^2\ge(G_1-1-a)^2\ge G_1^2-2(a+1)G_1$, where
$G_1-1-a=g+1-j_1-a\ge1$. Hence $(g+1-j)^22^{-j}>(G_1-2(a+1))/2^{a+1}=(g-j_1-2a)/2^{a+1}$, and dividing by $8(m+1)$ and using
$j_1\le\lceil\log_2(g+1)\rceil$ gives the bound.
\end{informalproof}

\paragraph{Why it matters.} Blindness bounds the visible fraction of a class from above by
$(g+3+\lceil\log_2(g+2)\rceil)/(m+1)$, and Theorem~V bounds it from below, for block labels, by about $g/(2^{a+4}(m+1))$:
both are of order $(m-k)/(m+1)$ up to logarithmic factors and constants depending on~$a$. So the coarse classes, a positive
fraction of all levels, see the other head's cancellations at full size on a positive fraction of their apex hops, on
average over shifts (see also \cref{rem:visible-numerics}). The next proposition says what this would give for block pairs,
but it needs more than Theorem~V provides. First, Conjecture~A$'$ fixes one shift for all classes at once. Second, the
proposition needs the cancellation in a window modulo~$q$, $-q/2<F\le-Q_h/2$, while Theorem~V gives only the integer event
$F\le-Q-1$, with no lower bound on~$F$: by the proof of \cref{lem:straddle}, $F=-Q-1+Y(o_u)-Y(o_p)$, and $Y$ sums
contributions of up to $m-\kappa$ nested levels, so $\abs F$ may grow linearly in~$m$, and the event is guaranteed to imply the
window only for $q$ of order $(m-\kappa)Q$ or larger.

\begin{remark}[Numerical: visible fractions]\label{rem:visible-numerics}
Measured fractions of apex hops with $F\le-Q-1$ were ten to a hundred times larger than the bound of Theorem~V.
\end{remark}

\begin{proposition}[Visibility at a fixed shift bounds the served levels of block pairs]\label{prop:Vfixed}
\notformalized
Let the heads carry $B(a_1,K_1)$ and $B(a_2,K_2)$ with $a_1,a_2\ge1$, at a fixed odd shift, on the $r$-fold cover with $r$
a multiple of~$q$, so that the two potentials define a label there modulo~$q$, and let $Q_1,Q_2$ be their block masses. Let $0<\beta<1$ and $\theta_0>0$, and suppose that for each head~$h'$ and each class
$(h',k)$ with $m-k\ge\beta(m+1)$, a fraction at least $c(\beta)>8\theta_0$ of its apex hops with $d<R_k$ have
$-q/2<F\le-Q_h/2$, where $h$ is the other head. Then at most $\beta(m+1)+2$ levels are doubly $\theta_0$-served.
\end{proposition}

\begin{informalproof}
Since $Q_h=(2^{a_h}-1)(2K_h+2)\ge2K_h+2$, the inequalities $K_1\ge Q_2/2$ and $K_2\ge Q_1/2$ cannot both hold: they would
give $Q_1\ge2K_1+2\ge Q_2+2$ and $Q_2\ge Q_1+2$. So some head~$h'$ has $K_{h'}<Q_h/2$ for the other head~$h$. Let $k\ge1$
with $m-k\ge\beta(m+1)$. By \cref{lem:band}(a) an apex hop with $d<R_k$ of a V of level~$k$ of head~$h'$ has error
$(A+F)\bmod q$ with $A\in\{0,K_{h'}\}$; on the hops of the hypothesis, $-q/2<A+F<0$, so the error lies in $[q/2,q)$. The
apex hops with $d<R_k$ are a fraction $(R_k-1)/(2R_k)\ge1/4$ of the class, so the class has total error at least
$c(\beta)\cdot\frac14\cdot\frac q2$ times its size, more than $\theta_0q$ times its size, and level~$k$ is not
$\theta_0$-served for head~$h'$. The levels $k\ge1$ with $m-k<\beta(m+1)$, together with level~$0$, number at most
$\beta(m+1)+2$.
\end{informalproof}

\subsection{A safe form of Conjecture A\texorpdfstring{$'$}{'}}\label{sec:dnp-App}

Blindness suggests, heuristically, that cheap labels may serve a number of fine levels that grows with~$m$; in the
measurements of \cref{rem:blind-numerics} the block pairs serve between $4$ and $8$ levels for $10\le m\le50$, above the
budget. The following form of the conjecture allows up to a fixed fraction of the levels to be served and keeps its
consequence.

\begin{conjecture}[Conjecture A$''$]\label{conj:App}
There are $\theta_0>0$, $\eta<1$ and constants $C,C'$ such that for all $s_0$ and~$m$ some odd shift~$\xi$ has the
following property: for all $r$ and~$q$, every label of cost at most $n_{\mathrm h}/8$ on the $r$-fold cover of the pair
$x_1=U(s_0,m,1)$, $x_2=x_1(\cdot+\xi)$, has at most $C\log_2q+C'+\eta(m+1)$ doubly $\theta_0$-served levels.
\end{conjecture}

If Conjecture~A$''$ holds and a label of cost at most $\eps n_{\mathrm h}$ with $\eps\le1/8$ and $\eps<(1-\eta)\theta_0/2$
exists on such a cover, then, as in \cref{sec:prelim-dnp}, it doubly serves more than $(m+1)(1-2\eps/\theta_0)$ levels, so
\[
  \log_2q\;>\;\frac{(1-\eta-2\eps/\theta_0)(m+1)-C'}{C}.
\]
By \cref{cor:wrap} this makes the pairs depth-robust as $m\to\infty$, and $\TCP(2)$ would fail; only the admissible
range of~$\eps$ shrinks, from $\theta_0/2$ to $(1-\eta)\theta_0/2$. \Cref{prop:Vfixed} proves Conjecture~A$''$ for pairs of
block labels, with $\eta=\beta+O(1/m)$ and no logarithmic term, from visibility in a window modulo~$q$ at a fixed shift.
Theorem~V does not supply this hypothesis, even on average over shifts: it gives the integer event $F\le-Q-1$ with no lower
bound on~$F$, hence the window only for $q$ of order $(m-\kappa)Q$, and only on average over shifts. Two steps remain open: a
bound on $\abs F$ on the straddling hops, or a version of Theorem~V that counts the window directly; and, at one fixed shift
for all coarse classes, equidistribution of a class's apex centres relative to the other head's canceller valleys of level
about $k+\log_2(m-k)$, a mixed-radix carry problem that we could not solve.

\subsection{The walk-DAG forms of the conjectures}\label{sec:dnp-wrap}

Conjectures~A$'$ and~A$''$ concern labels on the cyclic covers, where the chains of hops close up and Lemma~Zq is exact,
while $\TCP(2)$ concerns cuts of the walk DAG on $[0,T)$, where the chains are open. A label of the walk DAG costs at most
a constant more on the cover, so the conjectures already speak about the walk DAG; this is the elementary transfer in the
spirit of the fourth paper's remark on the cover and the walk DAG~\cite[Remark~5.23]{FourthPaper2026}, written for
served levels.

\begin{lemma}[From the walk DAG to the cover]\label{lem:wrap}
\notformalized
Let $s_0\ge m+2$, let $\xi$ be odd, let $N\ge1$ and $T=NP$, and let $\Psi\colon[0,T)\to\{0,\dots,q-1\}$ be a label of the
walk DAG of the pair $x_1=U(s_0,m,1)$, $x_2=x_1(\cdot+\xi)$ on $[0,T)$. Then $\Psi$, read as a label $G$ of the $N$-fold
cover, satisfies $\cost(G)\le\cost(\Psi)+2(2\Lambda+1)+1$.
\end{lemma}

\begin{informalproof}
Let $(p,u)$ be a hop of the cover, of head~$h$, lifted to~$\Z$ with $u\in[0,T)$, so that $p$ is the previous visit of
head~$h$ to the cell $x_h(u)$. If $p\ge0$, then $p$ is also the previous visit within $[0,T)$; since $U$ has no pauses,
$u\ge p+2$, so $(p,u)$ is a long hop of the walk DAG, and it is violated by~$G$ exactly when it is violated by~$\Psi$. If
$p<0$, then $u$ is the first visit of head~$h$ within $[0,T)$ to its cell; each head stays in the $2\Lambda+1$ cells of
$[-\Lambda,\Lambda]$, so there are at most $2(2\Lambda+1)$ such hops of the two heads. No pair is a hop of both heads, so
every long hop of the walk DAG is counted once on the cover: all turn times and turn positions are multiples of~$R_m$, so
a hop of~$x_1$ has $p+u\equiv0$ and a hop of~$x_2$ has $p+u\equiv-2\xi\pmod{R_m}$~\cite[Remark~5.23]{FourthPaper2026}, and
$2\xi\not\equiv0\pmod{R_m}$ because $\xi$ is odd and $R_m=2^{s_0-m}\ge4$. Finally, the descents of~$G$ at the times
$t\in[1,T)$ are those of~$\Psi$, and the cyclic comparison at $t=0$ adds at most one.
\end{informalproof}

\begin{corollary}[The conjectures on the walk DAG]\label{cor:wrap}
\notformalized
Suppose that Conjecture~A$''$ holds with constants $\theta_0$, $\eta$, $C$ and~$C'$. Let $0<\eps<\min(1/4,(1-\eta)\theta_0)$,
$\delta=((1-\eta)-\eps/\theta_0)/2$ and $N_0=\lceil1/\min(1/8-\eps/2,\delta\theta_0/2)\rceil$. For every~$m$, with $s_0=m+2$
and the shift~$\xi$ of the conjecture, and every $N\ge N_0$, every cut of at most $\eps NP$ edges of the walk DAG of the
pair on $[0,NP)$ has depth at least $2^{(\delta(m+1)-C')/C}-1$. So the deep-nesting pairs are depth-robust as $m\to\infty$
and $\TCP(2)$ fails. The same holds under Conjecture~A$'$, with $\eta=0$.
\end{corollary}

\begin{informalproof}
Let $F$ be a cut with at most $\eps NP$ edges and depth~$D$. By the converse of \cref{lem:labels}~\cite[Lemma~5.2]{FourthPaper2026}
there is a label~$\Psi$ with $q=D+1$ values and $\cost(\Psi)\le\eps NP$. By \cref{lem:wrap} the label~$G=\Psi$ of the
$N$-fold cover has $\cost(G)\le\eps NP+4\Lambda+3\le\eps_cn_{\mathrm h}$ with $\eps_c=\eps/2+1/(N(m+1))$, because
$n_{\mathrm h}=2NP$, $P=4(m+1)\Lambda$ and $4\Lambda+3\le8\Lambda$. Since $N\ge N_0$, $\eps_c\le1/8$ and
$2\eps_c/\theta_0\le\eps/\theta_0+\delta<1-\eta$, so $\eps_c<\theta_0/2$. By the per-class split
(\cref{sec:prelim-dnp}), $G$ doubly $\theta_0$-serves more than $(m+1)(1-2\eps_c/\theta_0)\ge(m+1)(1-\eps/\theta_0-\delta)$
levels, and Conjecture~A$''$ bounds their number by $C\log_2q+C'+\eta(m+1)$. As $1-\eta-\eps/\theta_0=2\delta$, this gives
$C\log_2q>\delta(m+1)-C'$, which is the claim.
\end{informalproof}

The proof uses Conjecture~A$''$ only through a bound on the number of doubly $\theta_0$-served levels by a function
of~$q$ plus $\eta(m+1)$. So a weaker statement already refutes $\TCP(2)$: that for every~$q$ there are $F(q)$
and~$m_0(q)$ such that for every $m\ge m_0(q)$, with $s_0=m+2$, some odd shift makes every label with $q$ values and cost at
most $n_{\mathrm h}/8$ on every cover doubly $\theta_0$-serve at most $F(q)+\eta(m+1)$ levels. Indeed, if $\TCP(2)$ held with
depth~$D$ at some $\eps<\min(1/4,(1-\eta)\theta_0)$, take $q=D+1$ and $m\ge\max(m_0(q),F(q)/\delta)$; the proof above gives
$F(q)+\eta(m+1)>(m+1)(\eta+\delta)$, a contradiction. The function~$F$ need not be monotone, and the shift may depend on~$q$.

